% !TEX program = xelatex
% Lernplatt - by Can Siebert
% Kurs 71 (Dig 42) - Tag 13 - Loesungsheft
% Nachhaltige IT-Services steuerbar machen: Service Portfolio, Capacity, Rightsizing
%
% Self-contained xelatex source. Kein biblatex, keine externen Bilder.

\documentclass[11pt,a4paper,oneside]{article}

\usepackage{polyglossia}
\setdefaultlanguage{german}

\usepackage[a4paper,margin=2.5cm]{geometry}

\usepackage{fontspec}
\defaultfontfeatures{Ligatures=TeX}

\IfFontExistsTF{Charter}{%
  \setmainfont{Charter}%
}{%
  \IfFontExistsTF{Noto Serif}{%
    \setmainfont{Noto Serif}%
  }{%
    \setmainfont{Liberation Serif}%
  }%
}

\IfFontExistsTF{Source Sans 3}{%
  \setsansfont{Source Sans 3}%
}{%
  \IfFontExistsTF{Source Sans Pro}{%
    \setsansfont{Source Sans Pro}%
  }{%
    \setsansfont{Liberation Sans}%
  }%
}

\newcommand{\caveatfontfallback}{\itshape}
\IfFontExistsTF{Caveat}{%
  \newfontfamily\caveatfont{Caveat}[Scale=1.15]%
}{%
  \newcommand\caveatfont{\caveatfontfallback}%
}

\setmonofont{Liberation Mono}

\usepackage{microtype}
\usepackage{eurosym}
\linespread{1.15}

\usepackage{xcolor}
\definecolor{lpInk}{RGB}{20,28,38}
\definecolor{lpAccent}{RGB}{157,74,26}
\definecolor{lpAccentLight}{RGB}{246,228,212}
\definecolor{lpAmber}{RGB}{222,138,30}
\definecolor{lpAmberLight}{RGB}{255,243,217}
\definecolor{lpAmberBorder}{RGB}{196,118,18}
\definecolor{lpGray}{RGB}{96,104,114}
\definecolor{lpGrayLight}{RGB}{238,240,243}
\definecolor{lpGreen}{RGB}{34,120,72}
\definecolor{lpGreenLight}{RGB}{222,238,228}
\definecolor{lpGreenBorder}{RGB}{24,96,58}

\definecolor{lpL1}{RGB}{34,120,72}
\definecolor{lpL2}{RGB}{22,90,140}
\definecolor{lpL3}{RGB}{170,42,52}

\usepackage[most]{tcolorbox}
\tcbuselibrary{breakable,skins}

\newtcolorbox{infobox}[1][Hinweis]{%
  enhanced, breakable,
  colback=lpAccentLight,
  colframe=lpAccent,
  boxrule=0.6pt,
  arc=2pt,
  left=10pt,right=10pt,top=8pt,bottom=8pt,
  fonttitle=\sffamily\bfseries\color{white},
  coltitle=white,
  title={#1},
  attach boxed title to top left={xshift=8pt,yshift=-8pt},
  boxed title style={size=small, colback=lpAccent, colframe=lpAccent, boxrule=0.4pt, arc=1pt},
  top=14pt
}

\newtcolorbox{loesungbox}[1][Musterl\"osung]{%
  enhanced, breakable,
  colback=lpGreenLight,
  colframe=lpGreenBorder,
  boxrule=0.6pt,
  arc=2pt,
  left=10pt,right=10pt,top=8pt,bottom=8pt,
  fonttitle=\sffamily\bfseries\color{white},
  coltitle=white,
  title={#1},
  attach boxed title to top left={xshift=8pt,yshift=-8pt},
  boxed title style={size=small, colback=lpGreenBorder, colframe=lpGreenBorder, boxrule=0.4pt, arc=1pt},
  top=14pt
}

\usepackage{enumitem}
\setlist{itemsep=2pt, topsep=4pt, parsep=0pt}
\setlist[itemize,1]{label={\textbullet},leftmargin=1.6em}
\setlist[itemize,2]{label={\textendash},leftmargin=1.4em}
\setlist[enumerate,1]{label=\arabic*., leftmargin=1.8em}

\usepackage{tabularx}
\usepackage{array}
\usepackage{booktabs}
\usepackage{longtable}

\usepackage{fancyhdr}
\usepackage{lastpage}
\pagestyle{fancy}
\fancyhf{}
\renewcommand{\headrulewidth}{0.3pt}
\renewcommand{\footrulewidth}{0pt}
\fancyhead[L]{\footnotesize\sffamily\color{lpGray}L\"osungsheft \textperiodcentered{} Kurs 71 \textperiodcentered{} Tag 13}
\fancyhead[R]{\footnotesize\sffamily\color{lpGray}Sustainable IT Services \& Capacity}
\fancyfoot[L]{\footnotesize\sffamily\color{lpGray}\textcopyright{} Can Siebert}
\fancyfoot[R]{\footnotesize\sffamily\color{lpGray}\thepage{} / \pageref{LastPage}}

\fancypagestyle{plain}{%
  \fancyhf{}%
  \renewcommand{\headrulewidth}{0pt}%
  \fancyfoot[C]{\footnotesize\sffamily\color{lpGray}\thepage{} / \pageref{LastPage}}%
}

\usepackage[explicit]{titlesec}

\titleformat{\section}[block]
  {\sffamily\Large\bfseries\color{lpAccent}}
  {\thesection.}{0.6em}{#1}
  [{\vspace{2pt}\color{lpAccent}\titlerule[0.6pt]}]

\titleformat{\subsection}[block]
  {\sffamily\large\bfseries\color{lpInk}}
  {\thesubsection}{0.6em}{#1}

\titlespacing*{\section}{0pt}{14pt}{8pt}
\titlespacing*{\subsection}{0pt}{10pt}{4pt}

\usepackage[hidelinks,unicode]{hyperref}

\newcommand{\akzent}[1]{{\caveatfont\color{lpAmber}#1}}
\newcommand{\fachbegriff}[1]{\textbf{#1}}

\newcommand{\niveaubadge}[2]{%
  \tcbox[on line, colback=#1, colframe=#1, coltext=white,
         boxrule=0pt, arc=2pt, left=5pt,right=5pt,top=1pt,bottom=1pt,
         nobeforeafter]{\sffamily\bfseries\footnotesize #2}%
}
\newcommand{\nivLi}{\niveaubadge{lpL1}{L1\,\textbullet\,Wissen}}
\newcommand{\nivLii}{\niveaubadge{lpL2}{L2\,\textbullet\,Anwendung}}
\newcommand{\nivLiii}{\niveaubadge{lpL3}{L3\,\textbullet\,Management}}

\newcommand{\zeit}[1]{%
  \tcbox[on line, colback=lpGrayLight, colframe=lpGrayLight, coltext=lpInk,
         boxrule=0pt, arc=2pt, left=5pt,right=5pt,top=1pt,bottom=1pt,
         nobeforeafter]{\sffamily\footnotesize\textbf{$\sim$\,#1}}%
}

\newcommand{\aufgabenkopf}[4]{%
  \par\vspace{6pt}
  \noindent{\sffamily\Large\bfseries\color{lpAccent}Aufgabe #1 \textendash{} #2}\par
  \vspace{2pt}
  \noindent{\color{lpAccent}\rule{\linewidth}{0.6pt}}\par
  \vspace{4pt}
  \noindent #3 \hfill \zeit{#4}\par
  \vspace{6pt}
}

\newcommand{\ytquery}[1]{%
  \par\vspace{4pt}
  \noindent{\footnotesize\sffamily\color{lpGray}\textbf{YouTube-Suchquery:} \texttt{#1}}\par
}

% =========================================================================
\begin{document}

\begin{titlepage}
  \pagestyle{empty}
  \vspace*{1.5cm}
  \noindent{\sffamily\footnotesize\color{lpGray}LERNPLATT \textperiodcentered{} BY CAN SIEBERT}\\[2pt]
  \noindent{\color{lpAccent}\rule{\linewidth}{1.2pt}}\\[4pt]
  \noindent{\sffamily\footnotesize\color{lpGray}L\"OSUNGSHEFT \textperiodcentered{} MODUL 6 \textperiodcentered{} TAG 13 VON 16 \textperiodcentered{} KURS 71 (Dig 42) \textperiodcentered{} 16.\,Juni 2026}

  \vspace{3cm}
  {\sffamily\fontsize{30}{34}\selectfont\bfseries\color{lpInk}L\"osungsheft\\Tag 13\par}

  \vspace{0.7cm}
  {\sffamily\large\color{lpGray}Musterl\"osungen zu den elf Aufgaben\\des Aufgabenhefts \textendash{} Service Portfolio,\\Capacity und Sustainable IT.\par}

  \vspace{1.5cm}
  {\caveatfont\LARGE\color{lpGreen}Musterl\"osungen \textperiodcentered{} nicht die einzige Antwort}\par

  \vfill

  \noindent\begin{minipage}{0.55\linewidth}
    {\sffamily\footnotesize\color{lpGray}Hinweis:}\\
    {\sffamily\small\color{lpInk}Musterl\"osungen sind Diskussionsbasis,\\nicht die einzige korrekte L\"osung.}\\[10pt]
    {\sffamily\footnotesize\color{lpGray}Begleitend:}\\
    {\sffamily\small\color{lpInk}Aufgabenheft \& Lehrskript Tag 13}
  \end{minipage}\hfill
  \begin{minipage}{0.4\linewidth}
    \raggedleft
    {\sffamily\footnotesize\color{lpGray}Dozent:}\\
    {\sffamily\small\color{lpInk}Can Siebert}\\[10pt]
    {\sffamily\footnotesize\color{lpGray}Niveau-Mix:}\\
    {\sffamily\small\color{lpInk}3\,L1 \textperiodcentered{} 5\,L2 \textperiodcentered{} 3\,L3}
  \end{minipage}

  \vspace{0.5cm}
  \noindent{\color{lpAccent}\rule{\linewidth}{0.6pt}}
\end{titlepage}

\section*{Hinweise zu den Musterl\"osungen}
\thispagestyle{plain}

\noindent Die folgenden Musterl\"osungen sind \emph{eine} m\"ogliche, didaktisch nachvollziehbare L\"osung. Insbesondere auf L2 und L3 sind mehrere Begr\"undungswege legitim, solange sie:

\begin{itemize}
  \item ITIL-4-Begriffe (Portfolio, Capacity, SLO/SLI/SLA) sauber nutzen,
  \item Annahmen explizit machen,
  \item Zielkonflikte (Verf\"ugbarkeit vs.\ Kosten vs.\ Nachhaltigkeit) benennen,
  \item zu einer klar formulierten Entscheidung mit Fr\"uhwarnsignal f\"uhren.
\end{itemize}

\begin{infobox}[Nutzung im Coaching]
Vergleichen Sie Ihre eigene Antwort \emph{vor} der Musterl\"osung. Wo weicht Ihre Argumentation ab? Ist die Abweichung eine andere Annahme \textendash{} oder ein Denkfehler? Genau dort entsteht der Lerneffekt.
\end{infobox}

\clearpage

% =========================================================================
% L1 AUFGABEN
% =========================================================================
\section*{Niveau L1 \textendash{} Wissen \& Reproduktion}
\addcontentsline{toc}{section}{L1 -- Wissen}

% --- AUFGABE 1 -----------------------------------------------------------
% LERNPLATT_HILFSVIDEOS
\section*{Hilfsvideos zu diesem Tag}
\addcontentsline{toc}{section}{Hilfsvideos zu diesem Tag}
\thispagestyle{plain}

\noindent Die folgenden YouTube-Erkl\"arvideos vermitteln die Inhalte, die Sie zur Bearbeitung der Aufgaben ben\"otigen. Klicken Sie auf den Titel, um das Video direkt zu \"offnen; die vollst\"andige URL steht in Klammern.

\begin{description}[style=nextline,leftmargin=18pt,labelindent=0pt,itemsep=8pt]
  \item[1.~ITIL 4 Service Portfolio Management — Pipeline, Live, Retired] \href{https://youtu.be/778SxGGGbrI}{\textcolor{lpAccent}{https://youtu.be/778SxGGGbrI}}\\[2pt]
  {\footnotesize\color{lpGray}(URL: \nolinkurl{https://youtu.be/778SxGGGbrI})}
  \item[2.~Nachhaltige Rechenzentren — Best Practices für mehr Energieeffizienz] \href{https://youtu.be/mW8RS3gND7o}{\textcolor{lpAccent}{https://youtu.be/mW8RS3gND7o}}\\[2pt]
  {\footnotesize\color{lpGray}(URL: \nolinkurl{https://youtu.be/mW8RS3gND7o})}
  \item[3.~ITIL Capacity Management — Praxis und Kennzahlen] \href{https://youtu.be/nP4Dk68pulM}{\textcolor{lpAccent}{https://youtu.be/nP4Dk68pulM}}\\[2pt]
  {\footnotesize\color{lpGray}(URL: \nolinkurl{https://youtu.be/nP4Dk68pulM})}
\end{description}
\vspace{0.6em}
\noindent{\footnotesize\color{lpGray}\textit{Tipp:} \"Offnen Sie das Video, lassen Sie es im Hintergrund laufen und arbeiten Sie parallel an der Aufgabe.}

\clearpage

\aufgabenkopf{1}{Service-Portfolio: Pipeline / Live / Retired}{\nivLi}{8\,min}

\ytquery{ITIL 4 service portfolio management pipeline live retired catalog}

\begin{loesungbox}
\textbf{1. Abgrenzung:}
\begin{itemize}
  \item \emph{Service Portfolio} \textendash{} die \emph{Gesamtmenge} aller Services (geplant, im Betrieb, abgek\"undigt). Steuerungssicht der IT-Leitung; enth\"alt auch interne Kosten, Risiken, Nachhaltigkeits\-treiber.
  \item \emph{Service Catalog} \textendash{} nur die aktuell \emph{verf\"ugbaren} (Live) Services. Kunden- und Bestellsicht; was kann ich heute beziehen?
  \item \emph{Service Pipeline} \textendash{} \emph{Vorhaben} und in Entwicklung befindliche Services; Planungssicht.
\end{itemize}

\smallskip
\textbf{2. Acht Mindestattribute:}

\smallskip
\begin{tabularx}{\linewidth}{@{}p{4.0cm}X@{}}
\toprule
\textbf{Attribut} & \textbf{Bedeutung} \\
\midrule
Service-Name              & sprechende Bezeichnung mit Bereichs-Pr\"afix \\
Owner                     & verantwortliche Rolle (nicht Person) \\
Lebenszyklus-Status       & Pipeline / Live / Retired \\
Kritikalit\"at / SLO-Klasse & Gold / Silver / Bronze; mit gekoppelten SLO-Zielen \\
Kosten pro Monat          & Run + Change (FinOps-tauglich) \\
Energieverbrauch          & Wert + Qualit\"atslabel (gemessen / gesch\"atzt) \\
Risiko                    & Betriebs-/Compliance-/Reputations-Risiko \\
Nachhaltigkeitshebel      & konkret (z.\,B.\ Off-Hours, Rightsizing, Retire) \\
\bottomrule
\end{tabularx}

\smallskip
\textbf{3. Begr\"undung Nachhaltigkeit als eigenes Attribut:} Wenn der Nachhaltigkeitseffekt nur in ``Kosten'' verbleibt, dominieren kurzfristige Sparma{\ss}nahmen \textendash{} ein Service kann \emph{billig im Run} und gleichzeitig \emph{teuer in CO\textsubscript{2}} sein. Eigenes Attribut macht die Trade-offs sichtbar und entscheidbar.
\end{loesungbox}

\clearpage

% --- AUFGABE 2 -----------------------------------------------------------
\aufgabenkopf{2}{Capacity \& Rightsizing \textendash{} Begriffe sortieren}{\nivLi}{8\,min}

\ytquery{IT capacity planning rightsizing SLO SLI SLA overprovisioning guardrail}

\begin{loesungbox}
\textbf{1. Begriffstabelle:}

\smallskip
\begin{tabularx}{\linewidth}{@{}p{2.8cm}X@{}}
\toprule
\textbf{Begriff} & \textbf{Definition + Beispiel} \\
\midrule
Capacity       & F\"ahigkeit, Last und SLOs zu erf\"ullen (Compute, Storage, Netz, People). Bsp.: 16 VMs, 200 GB RAM, 2 SREs. \\
Rightsizing    & Ressourcen an den \emph{realen} Bedarf anpassen ohne SLO zu verletzen. Bsp.: 10 VMs $\rightarrow$ 6 VMs nach 90-Tage-Auslastungsanalyse. \\
SLO            & Internes Zielwert pro SLI (z.\,B.\ Verf\"ugbarkeit $\geq$\,99,9\,\%). \\
SLI            & gemessener Indikator (z.\,B.\ tats\"achliche Verf\"ugbarkeit pro Monat). \\
SLA            & vertragliche Zusage gegen\"uber Kunde / Fachbereich (oft schw\"acher als SLO). \\
Overprovisioning & Ressourcen deutlich \"uber Bedarf; CPU $<$\,20\,\% Durchschnitt typisch. \\
Guardrail      & automatischer Stopp/Rollback bei Schwellverletzung; Bsp.: ``CPU $>$\,85\,\% \"uber 15 Min $\rightarrow$ Rollback''. \\
\bottomrule
\end{tabularx}

\smallskip
\textbf{2. Risiko ohne SLO / ohne Monitoring:}
\begin{itemize}
  \item \emph{Ohne SLO} ist Rightsizing Gl\"ucksspiel \textendash{} man wei{\ss} nicht, was ``gut genug'' bedeutet, und reduziert blind.
  \item \emph{Ohne Monitoring} entstehen typischerweise stille SLO-Verletzungen: der Service ist schon kaputt, bevor jemand das Fehlen von Alarmen bemerkt.
\end{itemize}

\smallskip
\textbf{3. Drei Anzeichen f\"ur Overprovisioning:}
\begin{itemize}
  \item Durchschnittliche CPU-Auslastung $<$\,20\,\% \"uber 90 Tage.
  \item Durchschnittliche RAM-Auslastung $<$\,40\,\% bei stabilem Lastprofil.
  \item Kein einziger Kapazit\"atsalarm in 90 Tagen \textendash{} obwohl Schwellen sauber konfiguriert sind.
\end{itemize}
\end{loesungbox}

\clearpage

% --- AUFGABE 3 -----------------------------------------------------------
\aufgabenkopf{3}{Acht Hebel f\"ur Energy-Efficient Data Centers}{\nivLi}{10\,min}

\ytquery{energy efficient data center virtualization rightsizing workload shift cooling}

\begin{loesungbox}
\textbf{1. Hebel mit Aufwand:}

\smallskip
\begin{tabularx}{\linewidth}{@{}p{4.0cm}p{1.4cm}X@{}}
\toprule
\textbf{Hebel} & \textbf{Aufwand} & \textbf{Was bedeutet das?} \\
\midrule
Konsolidierung         & Mid    & mehrere Workloads auf einer Plattform b\"undeln; weniger physische Maschinen. \\
Virtualisierung / Container & Mid & VMs / Container statt dedizierter Hardware; bessere Auslastung. \\
Abschalten Idle        & Low    & Non-Prod nachts und am Wochenende aus; trivial, aber selten konsequent. \\
K\"uhlung optimieren   & High   & Luftf\"uhrung, Hot/Cold-Aisle, Temperatur-Setpoint; baut tief in den RZ-Betrieb ein. \\
Hardware-Effizienzklassen& Mid  & beim n\"achsten Refresh effizientere Komponenten; einmalig wirksam. \\
Workload-Shift / Zeitfenster& Mid & Batch- und Backup-Last in Zeiten g\"unstigen Strommix legen. \\
Standort / Strommix    & High   & RZ-Standort bzw.\ Energievertrag mit Gr\"unanteil; mittelfristige Entscheidung. \\
Prozessdisziplin       & Low    & Change / Release / Decommissioning sauber f\"uhren \textendash{} keine ``Zombie-Server''. \\
\bottomrule
\end{tabularx}

\smallskip
\textbf{2. Drei Low Regret Moves:} \emph{Abschalten Idle} (Non-Prod off), \emph{Prozessdisziplin / Zombies dekommissionieren}, \emph{Workload-Shift / Batch-Fenster verschieben}. Hoher Effekt, niedriges Risiko, niedrige Investition.

\smallskip
\textbf{3. Zwei trugschl\"ussige ``Schnell-Hebel'':}
\begin{itemize}
  \item ``Alles in die Cloud schieben'' ohne Workload-Profil \textendash{} verschiebt das Problem oft nur, erh\"oht Kosten und kann CO\textsubscript{2}-Wirkung sogar verschlechtern, wenn der Cloud-Anbieter nicht in einem gr\"unen Strommix steht und Auto-Scaling fehlt.
  \item ``K\"uhlung runterfahren'' ohne Temperaturmonitoring \textendash{} kann Hardware-Ausf\"alle ausl\"osen; die SLO-Bru\-ch-Kosten \"ubersteigen das eingesparte kWh leicht um Gr\"o{\ss}enordnungen.
\end{itemize}
\end{loesungbox}

\clearpage

% =========================================================================
% L2 AUFGABEN
% =========================================================================
\section*{Niveau L2 \textendash{} Anwendung \& Transfer}
\addcontentsline{toc}{section}{L2 -- Anwendung}

% --- AUFGABE 4 -----------------------------------------------------------
\aufgabenkopf{4}{Service-Portfolio Canvas f\"ur 20 Services}{\nivLii}{25\,min}

\ytquery{IT service portfolio canvas critical SLO sustainability scoring}

\begin{loesungbox}
\textbf{Portfolio-Canvas (Auszug, acht Services):}

\smallskip
\begin{tabularx}{\linewidth}{@{}p{2.8cm}p{1.9cm}p{1.0cm}p{0.6cm}p{0.9cm}p{1.0cm}p{1.4cm}X@{}}
\toprule
\textbf{Service} & \textbf{Owner} & \textbf{Status} & \textbf{Krit.} & \textbf{SLO} & \textbf{Kosten} & \textbf{Energie} & \textbf{Entscheid.} \\
\midrule
Kundenportal      & Dig.\,Channel  & Live    & 5 & Gold   & 18\,k\,\euro & Compute       & Optimieren \\
Reporting/DWH     & FinOps         & Live    & 3 & Silver & 12\,k       & Batch-Last    & Optimieren \\
CRM               & Vertrieb       & Live    & 4 & Silver & 14\,k       & Compute       & Halten \\
Schaden-Mgmt      & Operations     & Live    & 5 & Gold   & 20\,k       & Compute+Storage & Halten \\
Dev/Test-Umg.     & Engineering    & Live    & 2 & Bronze & 8\,k        & Off-Hours-Hebel & Optimieren \\
Mobile App v2     & Dig.\,Channel  & Pipeline& 4 & Silver & 0           & Architekturwahl & Invest \\
Alt-BI            & FinOps         & Live    & 2 & Bronze & 5\,k        & Wenig Nutzung & Retire \\
Active Directory  & Security       & Live    & 5 & Gold   & 4\,k        & gering        & Halten \\
\bottomrule
\end{tabularx}

\smallskip
\textbf{Drei Entscheidungsempfehlungen mit Begr\"undung:}
\begin{itemize}
  \item \emph{Dev/Test-Umgebung optimieren:} Konsequenter Off-Hours-Shutdown + Image-Standardisierung; gesch\"atzt 35\,\% Energieersparnis ohne Auswirkung auf Engineering-Geschwindigkeit (Wiederanlauf $<$\,10 Min).
  \item \emph{Alt-BI retire:} Nutzung $<$\,5 User/Quartal, vollst\"andig durch Reporting/DWH abgedeckt; Stilllegung in 8 Wochen mit Migrationspfad und Archiv f\"ur Audit.
  \item \emph{Reporting/DWH optimieren:} Batch-Last in nachfragearme Zeitfenster verschieben + Storage-Tiering (\emph{Hot/Warm/Cold}); 20\,\% Kosten- und CO\textsubscript{2}-Hebel ohne SLO-Bruch.
\end{itemize}

\smallskip
\textbf{Zwei Zielkonflikte (durch Canvas sichtbar):}
\begin{itemize}
  \item \emph{Kritikalit\"at vs.\ Kosten:} Schaden-Mgmt (Krit.\,5) ist teuer, aber nicht verhandelbar.
  \item \emph{Verf\"ugbarkeit vs.\ Energie:} Gold-SLO erzwingt Redundanz \textendash{} R\"uckkopplung auf Energieverbrauch; sinnvolle Aufweichung nur bei Bronze-Services denkbar.
\end{itemize}
\end{loesungbox}

\clearpage

% --- AUFGABE 5 -----------------------------------------------------------
\aufgabenkopf{5}{Rightsizing-Plan ``Service auf 10 VMs''}{\nivLii}{25\,min}

\ytquery{rightsizing pilot rollout guardrails VM cpu memory peak SLO}

\begin{loesungbox}
\textbf{1. Drei Annahmen:}
\begin{itemize}
  \item Peak-Profil bleibt in den n\"achsten 90 Tagen vergleichbar (keine Marketing-Kampagne, kein Reporting-Stichtag\-wechsel).
  \item Autoscaling im Hypervisor m\"oglich (kann tempor\"ar VMs hochfahren bei Lastpeak).
  \item Lastprofil ist nicht von externen Triggern abh\"angig (kein synchroner Aufruf durch Drittsystem).
\end{itemize}

\smallskip
\textbf{2. Zweistufiger Rightsizing-Plan:}

\smallskip
\begin{tabularx}{\linewidth}{@{}p{2.2cm}p{4.0cm}p{2.6cm}X@{}}
\toprule
\textbf{Stufe} & \textbf{Aktion} & \textbf{Owner} & \textbf{Erfolg / Stopp} \\
\midrule
Pilot (Wo 1\,--\,3)   & 10 VMs $\rightarrow$ 7 VMs (3 stillgelegt); kleinere VM-Typen f\"ur 2 VMs; nachts nur 4 VMs aktiv & Service Owner + Plattformbetrieb & Erfolg: SLO eingehalten + Kosten\,$-$\,18\,\%; Stopp: SLO-Bruch oder Fehlerquote\,$>$\,0,5\,\% \\
Rollout (Wo 4\,--\,8) & wenn Pilot ok: Ausweitung auf 2 weitere \"ahnliche Services & Service Owner & Erfolg: gemittelte Energieersparnis\,$\geq$\,20\,\%; Stopp: Service-Owner-Veto bei einer Incident-Eskalation \\
\bottomrule
\end{tabularx}

\smallskip
\textbf{3. Zwei Guardrails:}
\begin{itemize}
  \item \emph{Antwortzeit-Guardrail:} 95.\ Perzentil der Antwortzeit $>$\,2\,s f\"ur l\"anger als 10 Minuten $\rightarrow$ automatischer Rollback (VMs wieder hochfahren), Ticket an Service Owner.
  \item \emph{CPU-Guardrail:} CPU \"uber alle aktiven VMs $>$\,85\,\% f\"ur mehr als 15 Minuten $\rightarrow$ Autoscale-Notfall + Alarm an Plattformbetrieb.
\end{itemize}

\smallskip
\textbf{4. Entscheidung unter Unsicherheit:}\\
\emph{Gr\"o{\ss}tes Risiko:} Peak-Profil unterz\"ahlt (z.\,B.\ Monatsende-Reporting taucht in den 90-Tage-Daten nicht ausreichend repr\"asentiert auf). \emph{Minimal-invasive Kompensation:}
\begin{itemize}
  \item Pilot \emph{startet erst nach} einem vollst\"andigen Monatsende-Zyklus, damit das Peak-Verhalten bestens dokumentiert ist.
  \item Notfall-Kapazit\"at (3 zus\"atzliche VMs) im Hypervisor \emph{vorgehalten} und in $<$\,10 Min aktivierbar; Plattformbetrieb hat dokumentierten Wiederanlauf.
\end{itemize}
\end{loesungbox}

\clearpage

% --- AUFGABE 6 -----------------------------------------------------------
\aufgabenkopf{6}{Sechs Ma{\ss}nahmen Capacity \& Rightsizing}{\nivLii}{25\,min}

\ytquery{ITIL capacity management sustainable IT rightsizing measures non-prod shutdown}

\begin{loesungbox}
\textbf{Sechs Ma{\ss}nahmen mit Owner / Guardrail / Monitoring / Rollback:}

\smallskip
\begin{tabularx}{\linewidth}{@{}p{3.0cm}p{2.4cm}p{3.0cm}p{2.6cm}X@{}}
\toprule
\textbf{Ma{\ss}nahme} & \textbf{Owner} & \textbf{Guardrail} & \textbf{Monitoring} & \textbf{Rollback} \\
\midrule
Non-Prod Off-Hours-Shutdown   & Plattformbetrieb   & Wiederanlauf $<$\,30\,Min m\"oglich & Off-Hours-Compliance-Quote & Wiederanlauf-Skript, on-demand \\
Rightsizing Pilot 2 Services  & Service Owner      & CPU $>$\,85\,\%, Fehlerquote $>$\,0,5\,\% & Synthetic Tests, p95-Latenz & 1-Klick-Restore VM-Snapshot \\
Storage-Tiering DWH           & FinOps + Storage   & Read-Latenz Hot $<$\,50\,ms      & Tier-Hit-Rate              & R\"uckverschiebung in Hot \\
Batch-Window verschieben      & Operations         & Berichts-Liefer-SLA halten       & Lieferzeitpunkt Reports     & alte Cron-Zeit reaktivieren \\
Retire Alt-BI                 & Portfolio Manager  & Alt-Reports archiviert, Audit\-Pfad& Anzahl aktive Nutzer        & Lesezugriff 90 Tage erhalten \\
Capacity-Review-Board (2-w\"ochentlich) & ITSM Lead & Entscheidungsprotokoll Pflichtfeld & Board-Stundenz\"ahler      & --- (governance) \\
\bottomrule
\end{tabularx}

\smallskip
\textbf{Priorisierung (Hebel vs.\ Aufwand):}
\begin{itemize}
  \item \emph{Wochen 1\,--\,4 (Quick Wins):} Non-Prod Shutdown, Capacity-Review-Board.
  \item \emph{Wochen 5\,--\,12 (Skalierung):} Rightsizing Pilot 2 Services, Batch-Window-Shift.
  \item \emph{Sp\"ater (Wo 13+):} Storage-Tiering DWH, Retire Alt-BI.
\end{itemize}
\emph{Begr\"undung:} Quick Wins liefern in 4 Wochen messbaren Effekt; Tiering und Retire haben Abh\"angigkeiten (Stakeholder-Kommunikation), daher sp\"ater.
\end{loesungbox}

\clearpage

% --- AUFGABE 7 -----------------------------------------------------------
\aufgabenkopf{7}{KPI-Set: 1 SLO, 1 Kosten, 1 Nachhaltigkeit}{\nivLii}{20\,min}

\ytquery{sustainable IT KPI green IT anti gaming coupling SLO cost energy}

\begin{loesungbox}
\textbf{1. Drei KPIs (kompakt):}

\smallskip
\begin{tabularx}{\linewidth}{@{}p{3.4cm}XX@{}}
\toprule
\textbf{KPI} & \textbf{Definition} & \textbf{Ziel / Owner / Review} \\
\midrule
SLO-Verf\"ugbarkeit (kritische Services) & Anteil der Stunden im Monat, in denen SLO eingehalten wurde & $\geq$\,99,9\,\% Gold-Services; Owner: SRE Lead; monatlich \\
Kosten pro Service / Monat        & FinOps-Run-Kosten / aktive Service-Instanz & Trend $-$\,15\,\% YoY; Owner: FinOps; monatlich \\
Non-Prod Off-Hours-Compliance     & \% der Stunden au{\ss}erhalb Kernzeit, in denen Non-Prod tats\"achlich aus war & $\geq$\,80\,\%; Owner: Plattformbetrieb; w\"ochentlich \\
\bottomrule
\end{tabularx}

\smallskip
\textbf{2. Kopplungsregel (Anti-Gaming):} \emph{Kostenreduktion z\"ahlt nur, wenn SLO-Verf\"ugbarkeit im selben Quartal $\geq$\,99,9\,\%.} Bonus / Reporting nur an gekoppeltes Tripel (Kosten $-$ SLO $-$ Off-Hours). Wer Kosten dr\"uckt und SLO bricht, hat per Definition \emph{nicht geliefert}.

\smallskip
\textbf{3. Greenwashing-Anti-Muster + Gegenma{\ss}nahme:}
\begin{itemize}
  \item \emph{Anti-Muster:} ``Non-Prod ist nachts aus'' \textendash{} aber das Team f\"ahrt es \emph{vor} dem n\"achsten Werktag fr\"uh wieder hoch, sodass Off-Hours-Wert k\"unstlich gut wird. Wirkung null.
  \item \emph{Gegenma{\ss}nahme:} Off-Hours-Z\"ahlung erfolgt aus Energie-/CPU-Daten, \emph{nicht} aus Tag/Nacht-Markern. Zus\"atzlich: \emph{tats\"achlich verbrauchte} kWh in Off-Hours als Z\"ahler.
\end{itemize}

\smallskip
\textbf{4. Vierte, optionale KPI:} \emph{Capacity-Decision Audit Coverage} \textendash{} \% der Capacity-/Rightsizing-Entscheidungen mit vollst\"andigem Entscheidungsprotokoll (Datenquelle, Verantwortlicher, Review-Datum). Ziel $\geq$\,90\,\%; sch\"utzt vor ``stillen'' Entscheidungen, die sp\"ater keine Begr\"undung haben.
\end{loesungbox}

\clearpage

% --- AUFGABE 8 -----------------------------------------------------------
\aufgabenkopf{8}{Fallstudie ``EuroHealth Services''}{\nivLii}{40\,min}

\ytquery{sustainable IT service portfolio capacity governance healthcare SLO energy}

\begin{loesungbox}
\textbf{1. Service-Portfolio (8 Services, kompakt):}

\smallskip
\begin{tabularx}{\linewidth}{@{}p{2.6cm}p{1.6cm}p{0.6cm}p{0.8cm}p{2.0cm}p{1.4cm}X@{}}
\toprule
\textbf{Service} & \textbf{Status} & \textbf{Krit.} & \textbf{SLO} & \textbf{Energie-Treiber} & \textbf{Risiko} & \textbf{Hebel} \\
\midrule
Patientenportal     & Live    & 5 & Gold   & Compute, Peak & Reputations   & Rightsizing + Caching \\
Klinikinformationssys.& Live  & 5 & Gold   & Compute       & Patientensafety & Konsolidierung \\
Reporting / DWH     & Live    & 3 & Silver & Batch-Last    & Audit-Verz\"og.& Batch-Shift, Tiering \\
DICOM-Archiv        & Live    & 4 & Silver & Storage       & Datenverlust  & Cold-Storage f\"ur Alt-Daten \\
Dev/Test            & Live    & 2 & Bronze & Compute       & gering        & Off-Hours-Shutdown \\
Mobile-App v2       & Pipeline& 4 & Silver & Architektur   & Time-to-Market& effiziente Defaults \\
Alt-HR-Tool         & Live    & 2 & Bronze & Compute       & Audit         & Retire 6 Monate \\
Active Directory    & Live    & 5 & Gold   & gering        & Sicherheits   & Halten \\
\bottomrule
\end{tabularx}

\smallskip
\textbf{2. Scoring (1\,--\,5; je h\"oher, desto besser/wichtiger). Top-3 zur Optimierung:} Reporting/DWH (Hebel 5, Risiko 2), Dev/Test (Hebel 5, Risiko 1), DICOM-Archiv (Hebel 4, Risiko 2). \textbf{Retire-Kandidat:} Alt-HR-Tool (Wert 1, Hebel 3, Risiko 2 \textendash{} kontrolliertes Auslaufen).

\smallskip
\textbf{3. Ma{\ss}nahmenpaket} (sechs Ma{\ss}nahmen \textendash{} siehe Aufgabe 6):
\begin{itemize}
  \item Non-Prod Off-Hours Shutdown (Dev/Test) \textendash{} Quick Win.
  \item Rightsizing Pilot Patientenportal (Off-Peak).
  \item Storage-Tiering DICOM-Archiv (Hot/Warm/Cold).
  \item Batch-Shift Reporting/DWH.
  \item Retire Alt-HR-Tool (Migrationspfad + Archiv).
  \item Capacity-Review-Board 2-w\"ochentlich (auditf\"ahig).
\end{itemize}

\smallskip
\textbf{4. Drei KPIs:}
\begin{itemize}
  \item \emph{SLO-Verf\"ugbarkeit Gold-Services} $\geq$\,99,95\,\% (Owner: SRE).
  \item \emph{Kosten pro Service} \textendash{} Trend $-$\,15\,\% in 4 Monaten (Owner: FinOps); gekoppelt: z\"ahlt nur bei SLO-Einhaltung.
  \item \emph{Off-Hours-Compliance Non-Prod} $\geq$\,80\,\% (Owner: Plattformbetrieb); gemessen aus kWh, nicht aus Schaltern.
\end{itemize}

\smallskip
\textbf{5. Entscheidung unter Unsicherheit:}
\begin{itemize}
  \item \emph{Sofort (4 Wochen):} Non-Prod Off-Hours-Shutdown + Capacity-Review-Board. Risiko niedrig, Hebel hoch, ohne Tooling-Abh\"angigkeit.
  \item \emph{Verschoben:} Retire Alt-HR-Tool (Stakeholder-Widerstand erwartet, ben\"otigt Migrationspfad) und Storage-Tiering DICOM (Datenintegrit\"at hat Vorrang \textendash{} ben\"otigt Validierungsphase).
\end{itemize}
\emph{Begr\"undung:} 15\,\%-Ziel ist in 4 Monaten realistisch erreichbar, wenn Quick Wins fr\"uh laufen und politisch heikle Aktionen sauber vorbereitet werden.
\end{loesungbox}

\clearpage

% =========================================================================
% L3 AUFGABEN
% =========================================================================
\section*{Niveau L3 \textendash{} Management \& Entscheidungsarchitektur}
\addcontentsline{toc}{section}{L3 -- Management}

% --- AUFGABE 9 -----------------------------------------------------------
\aufgabenkopf{9}{Portfolio Board \& Capacity Governance Design}{\nivLiii}{35\,min}

\ytquery{ITSM portfolio board capacity governance RACI decision protocol audit}

\begin{loesungbox}
\textbf{1. Portfolio Board Setup:}
\begin{itemize}
  \item \emph{Mitglieder:} Portfolio Owner (Vorsitz), Service-Owner-Vertretung (rotierend, 3 Bereiche), FinOps/Controlling, Operations/SRE Lead, Sustainable-IT Lead, Security Lead.
  \item \emph{Taktung:} monatlich 90 Min; bei Vorf\"allen ad hoc.
  \item \emph{Im Board entschieden:} Service-Aufnahme in Pipeline; SLO-Klassen-\"Anderung; Rightsizing-Ausnahmen; Retire-Empfehlungen.
  \item \emph{Eskaliert an Vorstand:} Budgetspr\"unge $>$\,250\,k\,\euro{}; SLO-Aufweichung Gold; Schliessung mit Kunden-/Compliance-Wirkung.
\end{itemize}

\smallskip
\textbf{2. RACI \"uber f\"unf Entscheidungstypen:}

\smallskip
\begin{tabularx}{\linewidth}{@{}p{3.4cm}ccccc@{}}
\toprule
 & \textbf{PortfOwn.} & \textbf{SvcOwn.} & \textbf{FinOps} & \textbf{Ops/SRE} & \textbf{Sec} \\
\midrule
Service einf\"uhren        & A    & R   & C   & C   & C \\
Service \"andern           & C    & R/A & C   & C   & C \\
Service stilllegen         & A    & R   & C   & I   & C \\
SLO-Klasse \"andern        & A    & R   & C   & C   & I \\
Rightsizing-Ausnahme       & I    & R/A & C   & C   & I \\
\bottomrule
\end{tabularx}

\smallskip
\textbf{3. Entscheidungsprotokoll (Pflichtfelder):}
\begin{itemize}
  \item \emph{Datum} und Teilnehmer (Name + Rolle).
  \item \emph{Entscheidung} im Klartext (was wurde beschlossen).
  \item \emph{Begr\"undung} + \emph{Datenquelle} (welcher Report, welche Metrik).
  \item \emph{Reviewdatum} (wann pr\"ufen wir, ob die Entscheidung wirkt).
  \item \emph{Verantwortlicher} f\"ur Umsetzung.
\end{itemize}

\smallskip
\textbf{4. Capacity Governance Bausteine:}
\begin{itemize}
  \item \emph{Demand Forecasting Minimal:} 12-Monats-Trend pro Service + Quartals-Adjust durch Service Owner. Kein KI-Modell n\"otig \textendash{} aber konsequente Datenpflege.
  \item \emph{Guardrail-Set:} CPU $>$\,85\,\% / Fehlerquote $>$\,0,5\,\% / p95-Latenz $>$\,2\,s; automatisch wirksam, dokumentiert im Service-Steckbrief.
  \item \emph{Ausnahmeprozess:} Service Owner kann Ausnahme f\"ur 30 Tage selbst genehmigen (mit Begr\"undung); l\"anger als 30 Tage nur durch Portfolio Board.
\end{itemize}

\smallskip
\textbf{5. Zwei Entscheidungen unter Unsicherheit:}
\begin{itemize}
  \item \emph{SLO-Reserve:} Gold-Services 30\,\% Headroom, Silver 15\,\%, Bronze 0\,\%. \emph{Annahme:} Peak-Streuung max.\ $+$25\,\% beobachtet; \emph{Fr\"uhwarnsignal:} mehr als zwei Headroom-Verletzungen pro Quartal $\rightarrow$ Headroom auf 35\,\% erh\"ohen, alternative Architektur (Edge-Caching) pr\"ufen. \emph{Eskalation:} CIO bei dritter Verletzung.
  \item \emph{Konsequente Optimierung trotz Widerstand:} Alt-HR-Tool und Alt-BI retire. \emph{Trade-off-Logik:} Nutzungswert $<$\,kritische Schwelle; alternative Wege existieren. \emph{Stakeholder-Plan:} 90-Tage-Lese-Zugriff, Migrationsworkshop, Champion aus dem betroffenen Bereich. \emph{Fr\"uhwarnsignal:} Widerstand eskaliert auf Bereichsleiter-Ebene $\rightarrow$ Sponsor-Gespr\"ach in $<$\,2 Wochen.
\end{itemize}
\end{loesungbox}

\clearpage

% --- AUFGABE 10 ----------------------------------------------------------
\aufgabenkopf{10}{ISO-20000-nahe Nachweislogik}{\nivLiii}{30\,min}

\ytquery{ISO 20000 ITSM audit evidence service management documentation compliance}

\begin{loesungbox}
\textbf{1. F\"unf Nachweisartefakte aus dem ITSM-Alltag:}

\smallskip
\begin{tabularx}{\linewidth}{@{}p{4.0cm}XX@{}}
\toprule
\textbf{Artefakt} & \textbf{Was gepr\"uft wird} & \textbf{Owner / Speicherort} \\
\midrule
Capacity-Board-Protokoll & Pflichtfelder vorhanden, Entscheidungen begr\"undet, Reviewdatum gesetzt & ITSM Lead / Confluence + Audit-Archiv \\
Change-Ticket mit Approver & Risikoklassifikation, Sign-off, Rollback-Plan, Post-Implementation-Review & Change Manager / ITSM-Tool \\
SLO-Report                 & SLI-Werte gemessen, Abweichungen erkl\"art, Massnahmen dokumentiert & SRE Lead / Monitoring-System \\
Rightsizing-Begr\"undung   & Datenbasis (90-Tage-Profil), Annahmen, Guardrails, Rollback & Service Owner / ITSM-Tool \\
Incident-Post-Mortem       & Zeitlinie, Root Cause, Massnahmen, Verantwortlicher, Follow-up-Tickets & Operations / Wiki \\
\bottomrule
\end{tabularx}

\smallskip
\textbf{2. Drei ``Papier-Anti-Muster'':}
\begin{itemize}
  \item \emph{Monatlicher Capacity-Report, den niemand liest:} 20 Seiten PDF an 30 Empf\"anger, ohne Frage / Entscheidung. \emph{Symptom:} keine R\"uckmeldung, keine Anschluss-Entscheidung.
  \item \emph{Doppelte Dokumentation in Wiki und Ticket:} dieselbe Information an zwei Stellen, einer veraltet \textendash{} und Auditor zieht beide.
  \item \emph{Manuell gepflegte Service-Liste:} Excel, drei verantwortliche Personen, niemand wei{\ss}, welche Version aktuell ist.
\end{itemize}

\smallskip
\textbf{3. Drei Automatik-Hebel:}
\begin{itemize}
  \item Ticket-System exportiert Change-Log direkt als Audit-Report (kein manuelles ``Zusammenstellen'').
  \item SLO-Dashboard wird auf Knopfdruck zum Audit-Report \textendash{} mit Signatur und Zeitstempel.
  \item Entscheidungs-Templates im ITSM-Tool zwingen Pflichtfelder vor Speicherung \textendash{} Capacity-Board-Protokoll entsteht ``nebenbei''.
\end{itemize}

\smallskip
\textbf{4. Faustregel:} Ein Nachweis ist \emph{Compliance}, wenn er aus einer \emph{Entscheidung} entstanden ist und einen Auditor in unter 2 Minuten zum Zweck f\"uhrt; er ist \emph{B\"urokratie}, sobald er die Entscheidung selbst nicht mehr beeinflusst und nur noch ``erzeugt'' wird, um ihn vorzeigen zu k\"onnen.
\end{loesungbox}

\clearpage

% --- AUFGABE 11 ----------------------------------------------------------
\aufgabenkopf{11}{Transferprojekt \textendash{} Sustainable IT Decision Architecture}{\nivLiii}{50\,min}

\ytquery{sustainable IT decision architecture CIO portfolio capacity governance roadmap}

\begin{loesungbox}
\emph{Musterl\"osung als Entscheidungsarchitektur; andere Konfigurationen sind m\"oglich.}

\smallskip
\textbf{1. Top-10 Entscheidungen:}
\begin{enumerate}
  \item Service-Priorit\"aten (Kritikalit\"at gekoppelt an Strategie, nicht Lautst\"arke).
  \item SLO-Klassen Gold/Silver/Bronze + Headroom-Regeln.
  \item Kapazit\"atsreserven pro Klasse (30\,\% / 15\,\% / 0\,\%).
  \item Rightsizing-Policy (Pflicht-Pilot vor Rollout, Guardrails verbindlich).
  \item Non-Prod-Regel (Off-Hours = Default; Ausnahmen mit Sunset).
  \item Invest vs.\ Retire (objektive Schwellen: Nutzer, Kosten, Risiko).
  \item Ausnahmeprozess (30 Tage Service Owner; l\"anger nur Board).
  \item Monitoring-Standards (SLI je Service Pflicht; sonst keine Prod-Freigabe).
  \item Risikoakzeptanz (Restrisiko sichtbar machen, nicht wegdefinieren).
  \item Review-Zyklen (monatlich Board, w\"ochentlich Ops, quartalsweise Vorstand).
\end{enumerate}

\smallskip
\textbf{2. Portfolio Governance:} Rollen wie Aufgabe 9; Board monatlich; Entscheidungs\-protokoll auditf\"ahig.

\smallskip
\textbf{3. Capacity \& Rightsizing Governance:} Demand Forecast Minimal; Guardrails verbindlich; Rollback-Plan vor Pilot; Eskalation an SRE + Service Owner.

\smallskip
\textbf{4. KPI \& Anti-Gaming:} SLO + Kosten + Nachhaltigkeit gekoppelt; Kostenreduktion nur g\"ultig bei SLO-Einhaltung; Off-Hours gemessen aus kWh, nicht aus Schaltern; Audit Coverage als Schutz-KPI.

\smallskip
\textbf{5. Roadmap:}
\begin{itemize}
  \item \emph{Wo 1\,--\,4 Quick Wins:} Non-Prod Off-Hours, Capacity Board konstituiert, KPI-Set freigegeben.
  \item \emph{Wo 5\,--\,12 Skalierung:} Rightsizing Pilot, Storage-Tiering, Retire-Welle 1.
  \item \emph{Sustain (laufend):} Monatliches Board, quartalsweiser Vorstandstrend, j\"ahrlicher Audit-Trockenlauf.
\end{itemize}

\smallskip
\textbf{6. Zwei Entscheidungen unter Unsicherheit:}
\begin{enumerate}
  \item \emph{SLO-Reserven (Headroom 30\,/\,15\,/\,0\,\%).} Annahme: Peak-Streuung max.\ $+$25\,\%. Verworfene Alternative: einheitliche 25\,\% (kostet zu viel). Fr\"uhwarnsignal: mehr als zwei Headroom-Verletzungen pro Quartal $\rightarrow$ Headroom anheben oder Architektur \"andern.
  \item \emph{Retire-Welle trotz Widerstand.} Annahme: $<$\,5 aktive Nutzer + alternative Wege vorhanden = legitim. Verworfene Alternative: ``alle behalten, weil sich jemand beschwert'' \textendash{} schiebt Kosten in Zukunft. Fr\"uhwarnsignal: Eskalation auf Bereichsleiter $\rightarrow$ Sponsor-Gespr\"ach + Stakeholder-Plan.
\end{enumerate}

\smallskip
\textbf{7. Anschluss an Strategie:}
\begin{itemize}
  \item Compliance: ISO-20000-nahe Nachweislogik schafft Audit-Sicherheit ohne B\"urokratie.
  \item Kosten: Hebel $-$15\,\% in 4 Monaten ist realistisch und vorzeigbar.
  \item Nachhaltigkeit: ESG-Reporting bekommt belastbare IT-Datenpunkte (CO\textsubscript{2}-Treiber pro Service), keine Greenwashing-Risiken.
\end{itemize}

\smallskip
\textbf{8. Vorstandssatz:}\\
``Wir f\"uhren ein gekoppeltes Steuerungssystem aus Service-Portfolio, Capacity Governance und SLO-/Kosten-/Nachhaltigkeits-KPIs ein, das in 12 Wochen 15\,\% Kosten \emph{ohne SLO-Bruch} freisetzt, Audit-Sicherheit schafft und Nachhaltigkeit messbar verankert \textendash{} Verantwortung ist je Service eindeutig, Restrisiken sind dokumentiert, und wir berichten quartalsweise an den Vorstand.''
\end{loesungbox}

\clearpage

\section*{Abschluss}
\thispagestyle{plain}

\noindent Die Musterl\"osungen sind eine Diskussionsbasis. Wenn Ihre eigene L\"osung anders ist, fragen Sie: \emph{Habe ich andere Annahmen \textendash{} oder eine andere Risikoabw\"agung?} Beides ist legitim, solange die Logik nachvollziehbar bleibt und Fr\"uhwarnsignale benannt sind.

\medskip
\begin{infobox}[N\"achster Schritt]
Bringen Sie Ihre L\"osungen in die Coaching-Sitzung. Leitfragen: Welche Annahmen liegen Ihren Klassifikationen zugrunde? Wo haben Sie bewusst Restrisiko akzeptiert? Welche Anti-Gaming-Logik sch\"utzt Ihr KPI-Set vor Schein-Erfolg?
\end{infobox}

\vspace{1cm}
\noindent{\color{lpAccent}\rule{\linewidth}{0.6pt}}\\[4pt]
{\footnotesize\sffamily\color{lpGray}Lernplatt \textperiodcentered{} by Can Siebert \textperiodcentered{} Kurs 71 (Dig 42) \textperiodcentered{} Tag 13 \textperiodcentered{} L\"osungsheft \textperiodcentered{} \textcopyright{} 2026}

\end{document}
